\chapter{Complexity — EXPCOM}
%%% generated by the blueprint pipeline from TCSlib.Complexity.Diagonalization.EXPCOM
\section{Overview}

Declarations of the Lean module \texttt{TCSlib.Complexity.Diagonalization.EXPCOM}.

\section{Declarations}

\begin{definition}[Uniform machine code]
\lean{Turing.UniformMachineCode}
\label{Turing.UniformMachineCode}
\leanok
\uses{Turing.CodeTM.toFinTM, Turing.EffectiveMachineCode, Turing.FinTM, Turing.FinTM.ComputesInTime, Turing.pairEncode}
A uniform machine code consists of an effective machine code, a Boolean finite-tape
machine, and a natural number $d$. For every pair of Boolean strings $\alpha,x$ and
every $t\in\mathbb N$, if the machine decoded from $\alpha$ outputs $[\mathrm{true}]$ on
$x$ in exactly $t$ steps, the simulator outputs $[\mathrm{true}]$ on the encoding of the
three inputs $\operatorname{bits}(t),\alpha,x$ in at most the specified time
$d\,(|\alpha|+|x|+t+1)^d$; if it does not output $[\mathrm{true}]$ on $x$ in exactly $t$
steps, the simulator outputs $[\mathrm{false}]$ on that encoding within the same time
bound.
\end{definition}

\begin{theorem}[Existence of a uniform machine code]
\lean{Turing.exists_uniformMachineCode}
\label{Turing.exists_uniformMachineCode}
\leanok
\uses{Turing.CodeTM.serialize, Turing.MachineFormat.exists_polyCanonizer, Turing.MachineFormat.exists_uniformSimulator, Turing.UniformMachineCode, Turing.codeDecode, Turing.codeDecode_serialize_pad, Turing.oneTapeMachineFormat}
\difficulty{4}
There exists a code for deterministic Boolean machines with one work tape, consisting of
an encoder and a total decoder such that decoding an encoded machine followed by any
number of $1$ bits returns that machine. The code also includes a finite-state Boolean
machine that computes the canonical serialization of each decoded machine within its
specified time bound, and a finite-state Boolean simulator with a natural number $d$.
For every program string $\alpha$, input string $x$, and time bound $t$, if the decoded
machine computes output $[1]$ on $x$ at exactly time $t$, the simulator computes $[1]$
on $\langle\langle\mathrm{bits}(t),\alpha\rangle,x\rangle$ at exactly
$d(|\alpha|+|x|+t+1)^d$ steps; if it does not compute $[1]$ at exactly time $t$, the
simulator computes $[0]$ on that input at exactly the same number of steps.
\end{theorem}

\begin{definition}[Chosen uniform machine code]
\lean{Complexity.expCode}
\label{Complexity.expCode}
\leanok
\uses{Turing.UniformMachineCode, Turing.exists_uniformMachineCode}
The object is a uniform machine code selected from the class of uniform machine codes
whose existence is asserted by the surrounding development.
\end{definition}

\begin{definition}[Exponential-time acceptance language]
\lean{Complexity.EXPCOM}
\label{Complexity.EXPCOM}
\leanok
\uses{Complexity.expCode, Turing.CodeTM.toFinTM, Turing.FinTM.ComputesInTime, Turing.pairEncode}
The language consists of exactly those Boolean strings $z$ for which there exist Boolean
strings $\alpha,x$ and a natural number $n$ such that $z$ is the pairing of $\alpha$
with the pairing of $x$ and the length-$n$ string of true bits, and the machine decoded
from $\alpha$ outputs $[\mathrm{true}]$ on $x$ in exactly $2^n$ steps.
\end{definition}

\begin{lemma}[Exponential padding bound]
\lean{Complexity.expcom_padding_bound}
\label{Complexity.expcom_padding_bound}
\leanok
\difficulty{5}
For all natural numbers $a,b,c,n$, $b\bigl(a\,2^{n^c}+1\bigr)^2\leq
2^{(b+2(a+1)+2)(n+1)^c}$.
\end{lemma}

\begin{theorem}[Exponential Time with an EXPCOM Oracle]
\lean{Complexity.EXP_subset_POracle_EXPCOM}
\label{Complexity.EXP_subset_POracle_EXPCOM}
\leanok
\uses{Complexity.EXP, Complexity.EXPCOM, Complexity.POracle, Complexity.PolyTimeComputable.pairEncode, Complexity.expCode, Complexity.expcom_padding_bound, Complexity.mem_POracle_of_polyTimeReducible, Complexity.polyTimeComputable_const, Complexity.polyTimeComputable_id, Complexity.polyTimeComputable_polyUnary, Turing.CodeTM.toFinTM, Turing.FinTM.ComputesInTime, Turing.FinTM.ComputesInTime.mono, Turing.FinTM.ComputesInTime.output_unique, Turing.FinTM.one_work_tape_binary, Turing.MachineCode.decode_encode, Turing.MultiTapeTM.indicator, Turing.exists_codeTM, Turing.pairEncode, Turing.pairEncode_injective}
\difficulty{6}
Every language decidable in time $2^{n^c}$ for some natural number $c$ is decidable in
polynomial time with oracle access to the following language. Fix the machine coding and
the machine code used in defining this language; a string belongs to the language
exactly when it has the form $\langle\alpha,\langle x,1^m\rangle\rangle$ and the machine
decoded from $\alpha$ halts on input $x$ after exactly $2^m$ steps with output
$[\mathrm{true}]$.
\end{theorem}

\begin{lemma}[Overflow characterizes the all-ones word]
\lean{Complexity.expcom_all_true}
\label{Complexity.expcom_all_true}
\leanok
\uses{Turing.incFixed, Turing.incFixed_eq_none_iff_replicate}
\difficulty{4}
Let $w$ be a finite Boolean word, and consider its fixed-width binary increment, with
overflow interpreted as the empty word. If the resulting word is empty, then every entry
of $w$ is true.
\end{lemma}

\begin{definition}[Second decoded component]
\lean{Complexity.expcom_tail}
\label{Complexity.expcom_tail}
\leanok
\uses{Complexity.pairSndD}
For a finite Boolean list $z$, define its tail by applying the defaulted
second-component operation twice to $z$.
\end{definition}

\begin{definition}[Validity predicate for encoded data]
\lean{Complexity.expcom_ok}
\label{Complexity.expcom_ok}
\leanok
\uses{Complexity.expcom_tail, Complexity.pairSndD, Turing.incFixed, Turing.pairDecode}
For a finite Boolean list $z$, define a Boolean predicate that holds exactly when $z$
and its defaulted second component both decode as pairs, and when applying the
fixed-increment operation to the tail obtained by two successive defaulted
second-component operations yields the empty list (with an empty list used if that
operation returns no value).
\end{definition}

\begin{definition}[Encoded simulation input]
\lean{Complexity.expcom_simIn}
\label{Complexity.expcom_simIn}
\leanok
\uses{Complexity.expcom_tail, Complexity.pairFstD, Complexity.pairSndD, Turing.incFixed, Turing.pairEncode}
For a finite Boolean list $z$, define an encoding obtained by pairing twice: the inner
pair consists of the fixed-increment transformation of the tail of $z$ with a final
false bit appended (using the empty list if the transformation returns no value), and
the defaulted first component of $z$; the outer pair consists of this inner encoding and
the defaulted first component of the tail of $z$.
\end{definition}

\begin{definition}[Validated simulation input]
\lean{Complexity.expcom_front}
\label{Complexity.expcom_front}
\leanok
\uses{Complexity.expcom_ok, Complexity.expcom_simIn, Turing.pairEncode}
For a finite Boolean list $z$, define its front to be its encoded simulation input when
the validity predicate holds, and otherwise the encoding obtained by pairing the empty
list with the empty list and then pairing that result with the empty list.
\end{definition}

\begin{lemma}[Characterization of valid front inputs]
\lean{Complexity.expcom_ok_iff}
\label{Complexity.expcom_ok_iff}
\leanok
\uses{Complexity.expcom_all_true, Complexity.expcom_ok, Complexity.expcom_tail, Complexity.pairSndD, Turing.eq_pairEncode_of_pairDecode, Turing.incFixed, Turing.incFixed_replicate_true_append, Turing.pairDecode, Turing.pairDecode_pairEncode, Turing.pairEncode}
\difficulty{5}
A Boolean word $z$ passes the front-input validity test if and only if there exist
Boolean words $\alpha,x$ and a natural number $n$ such that $z=\langle\alpha,\langle
x,1^n\rangle\rangle$.
\end{lemma}

\begin{lemma}[Polynomial-time computability of the front map]
\lean{Complexity.expcom_front_poly}
\label{Complexity.expcom_front_poly}
\leanok
\uses{Complexity.PolyTimeComputable, Complexity.PolyTimeComputable.comp, Complexity.PolyTimeComputable.pairEncode, Complexity.expcom_front, Complexity.expcom_ok, Complexity.expcom_simIn, Complexity.expcom_tail, Complexity.pairFstD, Complexity.pairFstD_pairEncode, Complexity.pairSndD_pairEncode, Complexity.polyTimeComputable_and, Complexity.polyTimeComputable_const, Complexity.polyTimeComputable_ite, Complexity.polyTimeComputable_of_linear, Complexity.polyTimeComputable_pairFstD, Complexity.polyTimeComputable_pairSndD, Complexity.polyTimeComputable_pairSwap, Turing.FinTM.computesFunInTime_appendBit, Turing.FinTM.computesFunInTime_ifEq, Turing.FinTM.computesFunInTime_incFixed, Turing.FinTM.computesFunInTime_pairValid, Turing.incFixed, Turing.pairDecode, Turing.pairEncode}
\difficulty{5}
The map that sends each Boolean word to its prescribed front transformation is
computable by a finite Turing machine in polynomial time: there exist a machine and
natural numbers $C,c$ such that, on every input of length $n$, it halts with output
equal to that transformation and uses exactly $C(n+1)^c$ steps.
\end{lemma}

\begin{lemma}[Front transformation on encoded triples]
\lean{Complexity.expcom_front_triple}
\label{Complexity.expcom_front_triple}
\leanok
\uses{Complexity.expcom_front, Complexity.expcom_ok, Complexity.expcom_ok_iff, Complexity.expcom_simIn, Complexity.expcom_tail, Turing.bits_two_pow_eq_replicate, Turing.incFixed, Turing.incFixed_replicate_true_append, Turing.pairEncode}
\difficulty{3}
For all Boolean words $\alpha,x$ and all $n\in\mathbb{N}$, applying the front
transformation to the encoding of $\alpha$, $x$, and a word of $n$ true bits gives the
encoding of the pair whose first component is the encoding of the binary representation
of $2^n$ together with $\alpha$, and whose second component is $x$; explicitly,
$\mathrm{front}(\langle\alpha,\langle
x,1^n\rangle\rangle)=\langle\langle\mathrm{bits}(2^n),\alpha\rangle,x\rangle$.
\end{lemma}

\begin{lemma}[Exponential bound for the time budget]
\lean{Complexity.expcom_decider_budget}
\label{Complexity.expcom_decider_budget}
\leanok
\difficulty{5}
For all $C,c,d,m\in\mathbb{N}$, $$2\bigl(C(m+1)^c\bigr)+d(m+2^m+1)^d+2\le
(2C+d+2)\,2^{(c+d)(m+1)}.$$
\end{lemma}

\begin{lemma}[Membership of Encoded Inputs in EXPCOM]
\lean{Complexity.expcom_mem_triple}
\label{Complexity.expcom_mem_triple}
\leanok
\uses{Complexity.EXPCOM, Complexity.expCode, Turing.CodeTM.toFinTM, Turing.FinTM.ComputesInTime, Turing.pairEncode, Turing.pairEncode_injective}
\difficulty{4}
For Boolean words $\alpha$ and $x$ and a natural number $n$, the word
$\langle\alpha,\langle x,1^n\rangle\rangle$ belongs to $\mathrm{EXPCOM}$ if and only if
the machine decoded from $\alpha$ computes the output word $[1]$ on input $x$ in exactly
$2^n$ steps.
\end{lemma}

\begin{lemma}[Simulator Runtime on Transformed Inputs]
\lean{Complexity.expcom_simulator_on_image}
\label{Complexity.expcom_simulator_on_image}
\leanok
\uses{Complexity.EXPCOM, Complexity.expCode, Complexity.expcom_front, Complexity.expcom_front_triple, Complexity.expcom_mem_triple, Complexity.expcom_ok, Complexity.expcom_ok_iff, Turing.CodeTM.toFinTM, Turing.FinTM.ComputesInTime, Turing.FinTM.ComputesInTime.mono, Turing.FinTM.not_computesInTime_zero, Turing.MultiTapeTM.indicator, Turing.length_pairEncode, Turing.pairEncode}
\difficulty{5}
For every Boolean word $z$, the fixed simulator computes, on the front transformation of
$z$, the one-bit output whose value is the indicator of whether $z$ belongs to
$\mathrm{EXPCOM}$. Its running time is exactly $d(|z|+2^{|z|}+1)^d$, where $d$ is the
simulator's degree parameter.
\end{lemma}

\begin{theorem}[Exponential-Time Decision of Bounded Acceptance]
\lean{Complexity.expcom_decider}
\label{Complexity.expcom_decider}
\leanok
\uses{Complexity.EXPCOM, Complexity.expCode, Complexity.expcom_decider_budget, Complexity.expcom_front, Complexity.expcom_front_poly, Complexity.expcom_simulator_on_image, Turing.FinTM, Turing.FinTM.ComputesInTime.mono, Turing.FinTM.DecidesInTime, Turing.FinTM.exists_comp_on_image, Turing.MultiTapeTM.indicator}
\difficulty{3}
Let $\mathrm{EXPCOM}$ be the language of Boolean words $z$ for which there are Boolean
words $\alpha,x$ and a natural number $n$ such that $z=\langle\alpha,\langle
x,1^n\rangle\rangle$ and the machine decoded from $\alpha$ accepts $x$ within $2^n$
steps. There exist a finite Boolean Turing machine $D$ and natural numbers $A,B$ such
that, for every input of length $m$, $D$ decides membership in $\mathrm{EXPCOM}$ within
$A\,2^{B(m+1)}$ steps.
\end{theorem}

\begin{theorem}[Constant functions are exponential-polynomial]
\lean{Complexity.expcom_expPoly_const}
\label{Complexity.expcom_expPoly_const}
\leanok
\uses{Complexity.ExpPoly, Complexity.ExpPoly.of_le, Complexity.expPoly_poly}
\difficulty{2}
For every $A\in\mathbb{N}$, the constant function $n\mapsto A$ is
exponential-polynomial: there exist $K,k\in\mathbb{N}$ such that $A\le 2^{K(n+1)^k}$ for
every $n\in\mathbb{N}$.
\end{theorem}

\begin{theorem}[Exponential-polynomial bound with a constant factor]
\lean{Complexity.expcom_expPoly_mul_two_pow}
\label{Complexity.expcom_expPoly_mul_two_pow}
\leanok
\uses{Complexity.ExpPoly}
\difficulty{4}
For all $A,K,k\in\mathbb{N}$, the function $n\mapsto A\,2^{K(n+1)^k}$ is
exponential-polynomial: there exist $K',k'\in\mathbb{N}$ such that $A\,2^{K(n+1)^k}\le
2^{K'(n+1)^{k'}}$ for every $n\in\mathbb{N}$.
\end{theorem}

\begin{theorem}[Exponential-polynomial summit bound]
\lean{Complexity.expcom_summit_budget}
\label{Complexity.expcom_summit_budget}
\leanok
\uses{Complexity.ExpPoly, Complexity.ExpPoly.add, Complexity.ExpPoly.mul, Complexity.ExpPoly.of_le, Complexity.expPoly_poly, Complexity.expcom_expPoly_const, Complexity.expcom_expPoly_mul_two_pow}
\difficulty{6}
For all natural numbers $a',A,B,b,e,b',C,k$, define
$T(n)=b'2^{C(n+1)^k}\Bigl(b\bigl((n+C(n+1)^k+1)^e+a'(2(n+C(n+1)^k)+3)\bigl(2(n+C(n+1)^k)+3+A2^{B(2(n+C(n+1)^k)+3)}\bigr)\bigr)+(n+C(n+1)^k+1)^{k+1}\Bigr)$.
Then $T$ is exponential-polynomial: there exist natural numbers $K,j$ such that
$T(n)\leq 2^{K(n+1)^j}$ for every natural number $n$.
\end{theorem}

\begin{theorem}[Nondeterministic oracle exponential-time inclusion]
\lean{Complexity.NPOracle_EXPCOM_subset_EXP}
\label{Complexity.NPOracle_EXPCOM_subset_EXP}
\leanok
\uses{Complexity.EXP, Complexity.EXPCOM, Complexity.ExpPoly.mem_EXP, Complexity.NPOracle, Complexity.exists_proj_decider, Complexity.expcom_decider, Complexity.expcom_summit_budget, Turing.Cfg.ofWords, Turing.FinOracleNDTM.choiceOut, Turing.FinOracleNDTM.exists_choiceTM, Turing.FinOracleNDTM.mem_iff_exists_choiceOut, Turing.FinOracleTM.exists_plain_simulation_of_monotone, Turing.FinTM.ComputesInTime.mono, Turing.FinTM.certificateShell_proj, Turing.FinTM.exists_certificateShell, Turing.FinTM.prefixWords_update, Turing.MultiTapeTM.runFrom, Turing.MultiTapeTM.runFrom_add, Turing.MultiTapeTM.runFrom_of_halt}
\difficulty{6}
Every language decidable by a nondeterministic Turing machine with oracle
$\mathrm{EXPCOM}$ in time at most $c(n^c+1)$ for some natural number $c$ belongs
to $\mathrm{EXP}$, the class of languages decidable by a deterministic Turing machine
in time at most $d\,2^{n^d}$ for some natural number $d$. Here $\mathrm{EXPCOM}$
consists of strings encoding a machine description, an input, and a unary time bound
$n$, such that the described machine accepts that input within $ 2^n$ steps.
\end{theorem}

\begin{theorem}[Polynomial oracle characterization of exponential time]
\lean{Complexity.POracle_EXPCOM_eq_EXP}
\label{Complexity.POracle_EXPCOM_eq_EXP}
\leanok
\uses{Complexity.EXP, Complexity.EXPCOM, Complexity.EXP_subset_POracle_EXPCOM, Complexity.NPOracle_EXPCOM_subset_EXP, Complexity.POracle, Complexity.POracle_subset_NPOracle}
\difficulty{1}
Let $\mathrm{EXPCOM}$ consist of strings encoding a machine description $\alpha$, an
input $x$, and a unary time bound $n$, such that the machine described by $\alpha$
accepts $x$ within $ 2^n$ steps. Then the class of languages decidable by
deterministic Turing machines with oracle $\mathrm{EXPCOM}$ in time at most
$c(n^c+1)$ for some natural number $c$ is exactly $\mathrm{EXP}$, the class of
languages decidable in deterministic time at most $d\,2^{n^d}$ for some natural number
$d$.
\end{theorem}

\begin{theorem}[Oracle nondeterminism characterizes exponential time]
\lean{Complexity.NPOracle_EXPCOM_eq_EXP}
\label{Complexity.NPOracle_EXPCOM_eq_EXP}
\leanok
\uses{Complexity.EXP, Complexity.EXPCOM, Complexity.NPOracle, Complexity.NPOracle_EXPCOM_subset_EXP, Complexity.POracle_EXPCOM_eq_EXP, Complexity.POracle_subset_NPOracle}
\difficulty{3}
The class of Boolean languages decidable by a nondeterministic oracle Turing machine
with oracle $\mathrm{EXPCOM}$, in time bounded by a constant multiple of $n^c+1$ for
some natural number $c$, is exactly $\mathrm{EXP}$, the union over natural numbers $c$
of deterministic time classes with bounds a constant multiple of $2^{n^c}$.
\end{theorem}

\begin{theorem}[Deterministic and nondeterministic exponential-oracle classes]
\lean{Complexity.POracle_EXPCOM_eq_NPOracle_EXPCOM}
\label{Complexity.POracle_EXPCOM_eq_NPOracle_EXPCOM}
\leanok
\uses{Complexity.EXPCOM, Complexity.NPOracle, Complexity.NPOracle_EXPCOM_eq_EXP, Complexity.POracle, Complexity.POracle_EXPCOM_eq_EXP}
\difficulty{1}
For the oracle language $\mathrm{EXPCOM}$, the class of languages decidable in
polynomial time by deterministic oracle Turing machines is equal to the class decidable
in polynomial time by nondeterministic oracle Turing machines:
$\mathrm{P}^{\mathrm{EXPCOM}}=\mathrm{NP}^{\mathrm{EXPCOM}}$.
\end{theorem}
